% !TeX program = xelatex
% SPDX-License-Identifier: MIT
% Copyright (c) 2026 Ziyu Peng and 刘毅涵
% Author: https://pzy54154631-hub.github.io/homepage/
% Author: https://pzy54154631-hub.github.io/liuyihan/
% 自设会计教学案例；金额单位为元，不对应任何企业。
\documentclass[UTF8,a4paper,fontset=fandol]{ctexart}
\usepackage[margin=2.2cm]{geometry}
\usepackage{amsmath,amssymb,booktabs,longtable,array}
\usepackage[hidelinks]{hyperref}
\title{会计分录与报表：让同一组数字相互勾稽}
\author{\href{https://pzy54154631-hub.github.io/homepage/}{\textit{Ziyu Peng}}, University of Colorado Boulder
  \\[0.35em]
  \href{https://pzy54154631-hub.github.io/liuyihan/}{刘毅涵}，暨南大学经济学院}
\date{}
\begin{document}
\maketitle
假设一家模拟服务工作室期初余额全部为零。下面记录一个教学期间的交易，
不考虑税费，期末无分红；设备本期折旧给定为 100。
耗材购入时先确认为资产，本期耗用部分才转为费用。
所有数值均为自设，报表是为解释勾稽关系而简化的教学版。

\section{分录：每笔借贷相等}
\begin{longtable}{@{}p{1cm}p{4.9cm}rr@{}}
\caption{模拟交易分录（单位：元）}\label{tab:entries}\\
\toprule
事项 & 科目 & 借方 & 贷方\\
\midrule
\endfirsthead
\toprule
事项 & 科目 & 借方 & 贷方\\
\midrule
\endhead
\midrule
\multicolumn{4}{r}{续下页}\\
\endfoot
\bottomrule
\endlastfoot
1 & 银行存款 & 10000 & \\
  & \quad 投入资本 & & 10000\\
\addlinespace
2 & 设备 & 3000 & \\
  & \quad 银行存款 & & 3000\\
\addlinespace
3 & 耗材 & 600 & \\
  & \quad 银行存款 & & 600\\
\addlinespace
4 & 耗材 & 500 & \\
  & \quad 应付账款 & & 500\\
\addlinespace
5 & 银行存款 & 1800 & \\
  & 应收账款 & 600 & \\
  & \quad 服务收入 & & 2400\\
\addlinespace
6 & 租金费用 & 500 & \\
  & \quad 银行存款 & & 500\\
\addlinespace
7 & 耗材费用 & 200 & \\
  & \quad 耗材 & & 200\\
\addlinespace
8 & 折旧费用 & 100 & \\
  & \quad 累计折旧 & & 100\\
\end{longtable}

\section{本期利润与期末权益}
\begin{center}
\begin{tabular}{lr}
\toprule
利润表项目 & 金额（元）\\
\midrule
服务收入 & 2400\\
减：租金费用 & 500\\
减：耗材费用 & 200\\
减：折旧费用 & 100\\
\midrule
本期利润 & 1600\\
\bottomrule
\end{tabular}
\end{center}
在本例无税费的设定下，本期利润为 1600。期末权益为
$0+10000+1600-0=11600$，依次对应期初权益、投入、本期利润和分配。

\section{期末资产负债表}
\begin{center}
\begin{tabular}{lr}
\toprule
资产项目 & 金额（元）\\
\midrule
银行存款 & 7700\\
应收账款 & 600\\
耗材 & 900\\
设备原值 & 3000\\
减：累计折旧 & 100\\
\midrule
资产总计 & 12100\\
\midrule
负债和权益项目 & 金额（元）\\
\midrule
应付账款 & 500\\
投入资本 & 10000\\
本期形成的留存利润 & 1600\\
\midrule
负债和权益合计 & 12100\\
\bottomrule
\end{tabular}
\end{center}
核对 $12100=500+11600$。累计折旧抵减设备账面价值，并非负债。

\clearpage
\section{本期现金流量汇总}
\begin{center}
\begin{tabular}{lr}
\toprule
项目 & 金额（元）\\
\midrule
经营：收到服务款 & 1800\\
经营：购入耗材支付 & -600\\
经营：支付租金 & -500\\
经营活动现金流量净额 & 700\\
\midrule
投资：购买设备 & -3000\\
筹资：收到投入资本 & 10000\\
\midrule
现金净增加额 & 7700\\
期初现金余额 & 0\\
期末现金余额 & 7700\\
\bottomrule
\end{tabular}
\end{center}
赊购耗材 500 不产生本期现金收付。折旧 100 也不直接产生现金收付。
本例银行存款即现金流量表中的现金，未另设现金等价物。
以间接思路核对经营现金流：
\[
1600+100-600-900+500=700.
\]
各项依次为利润、折旧、应收增加、耗材增加和应付增加。
\end{document}
